\section{Related Work}
\label{sec:related}

\paragraph{Texture segmentation and benchmarks.} Classical texture recognition and segmentation work established the core datasets and problem structure that still matter here, including DTD, MINC, and the RWTD/STLD benchmark lineage~\cite{cimpoi2014dtd,bell2015minc,cimpoi2016deepfilterbanks,khan2015stld,khan2017c2f,khan2018learnedstld,mikes2022texturebenchmark}. Much of that literature focuses either on handcrafted/local descriptors or on fully supervised dense predictors specialized to a fixed texture regime~\cite{andrearczyk2017texturefcn,liu2019bowtocnn}. Our question is different: given a frozen foundation segmenter, how much texture partition structure is already recoverable before any backbone retraining?

\paragraph{SAM-family adaptation.} Segment Anything made high-recall proposal banks widely available, and SAM 2 extended that foundation-segmentation interface to stronger image and video settings~\cite{kirillov2023sam,ravi2024sam2}. Most SAM adaptation work improves the foundation model itself through prompt tuning, adversarial tuning, or higher-quality decoders~\cite{ke2023hqsam,zhang2023persam,li2024asam}. TextureSAM is the most relevant prior route for this paper because it specifically injects texture behavior into the SAM backbone and establishes a strong baseline on the RWTD/STLD family~\cite{cohen2025texturesam,cohen2026texturesam}. We do not argue against that line. We ask whether recoverable texture behavior also lives above a frozen proposal engine.

\paragraph{Representation probing and proposal reasoning.} A small clustering-based probe can answer a different question from a task-trained decoder: not ``can a large head solve the benchmark,'' but ``what structure is already latent in the frozen representation?'' That perspective motivates our feature-space exploration route. Our proposal-/mask-space exploration route instead connects texture segmentation to a more general class of reasoning problems over candidate masks: compatibility scoring, conservative selection, and selective completion from a frozen proposal bank. In this paper those two views are unified rather than treated as unrelated methods. The feature route asks what texture structure is already present in the representation; the proposal route asks what texture structure has already been externalized into the masks that the frozen model emits.
